% BHU PYQ -- Auto-generated & error-corrected
\documentclass[11pt,a4paper]{article}

\usepackage[a4paper,top=2.5cm,bottom=2.5cm,left=2.8cm,right=2.8cm,headheight=14pt]{geometry}
\usepackage{amsmath,amssymb}
\usepackage[T1]{fontenc}
\usepackage[utf8]{inputenc}
\usepackage{lmodern}
\usepackage{microtype}
\usepackage{fancyhdr}
\usepackage{enumitem}
\usepackage{hyperref}
\usepackage{lastpage}
\usepackage{parskip}

\hypersetup{colorlinks=true,urlcolor=black,linkcolor=black,
  pdftitle={STA/STB-502: Statistical Inference and Decision Theory -- BHU PYQ},pdfauthor={Banaras Hindu University}}

\pagestyle{fancy}\fancyhf{}
\renewcommand{\headrulewidth}{0.4pt}\renewcommand{\footrulewidth}{0.4pt}
\fancyhead[L]{\small   \textit{Banaras Hindu University}}
\fancyhead[R]{\small   \textit{STA/STB-502: Statistical Inference \& Decision Theory}}
\fancyfoot[C]{\small Page  \thepage\ of \pageref{LastPage}}


\newlist{parts}{enumerate}{2}
\setlist[parts,1]{label=(\alph*),leftmargin=*,itemsep=5pt,topsep=3pt}
\setlist[parts,2]{label=(\roman*),leftmargin=*,itemsep=3pt,topsep=2pt}

\newcommand{\pts}[1]{\hfill   \textit{[#1]}}


\begin{document}


\begin{center}
  {\large\bfseries BANARAS HINDU UNIVERSITY}\\[2pt]
  {
\normalsize B.A./B.Sc. (Hons.) Semester V Examination, 2013-14}\\[6pt]
  \rule{0.8\linewidth}{0.5pt}\\[6pt]
  {\large\bfseries Paper No. STA/STB-502: Statistical Inference and Decision Theory}\\[4pt]
  {
\normalsize Time: 3 Hours\hfill Maximum Marks: 70}
\end{center}
\rule{\linewidth}{0.4pt}
\smallskip



\noindent   \textit{Instructions: Answer any   \textbf{five} questions. All questions carry equal marks. Symbols have their usual meanings.}
\bigskip

\medskip


\noindent   \textbf{Question 1.}

\begin{parts}
  \item Define an unbiased estimator. What are the criteria of a good estimator? Show that $2ar{X}$ is an unbiased estimator of $ \theta$ in the rectangular distribution $U(0,  \theta)$. \pts{7}
  \item State and prove the invariance property of a consistent estimator. Let $X_1, X_2, \ldots, X_n$ be a random sample from a Bernoulli population with parameter $ \theta$. Show that $ar{X}(1 - ar{X})$ is a consistent estimator of $ \theta(1- \theta)$. \pts{7}
\end{parts}

\medskip\rule{\linewidth}{0.2pt}

\medskip


\noindent   \textbf{Question 2.}

\begin{parts}
  \item Define a sufficient statistic. Let $X_1, X_2, \ldots, X_n$ be a random sample of size $n$ from a Poisson distribution with parameter $ \theta$. Show that $\sum_{i=1}^n X_i$ is a complete sufficient statistic for $ \theta$. \pts{7}
  \item Let $X_1, X_2, \ldots, X_n$ be a random sample from the probability density function:
  \[
    f(x,  \theta) = \frac{2x}{ \theta^2}, \quad 0 < x <  \theta, \quad  \theta > 0.
  \]
  Find the minimum variance unbiased estimator (MVUE) of $ \theta$. \pts{7}
\end{parts}

\medskip\rule{\linewidth}{0.2pt}

\medskip


\noindent   \textbf{Question 3.}

\begin{parts}
  \item Stating the regularity conditions, state and prove the Cram\'er--Rao inequality. \pts{7}
  \item Let $X_1, X_2, \ldots, X_n$ be a random sample from the probability density function:
  \[
    f(x,  \theta) =  \theta e^{- \theta x}, \quad x \geq 0, \quad  \theta > 0.
  \]
  Obtain the minimum variance bound (MVB) estimator of $ \theta$. \pts{7}
\end{parts}

\medskip\rule{\linewidth}{0.2pt}

\medskip


\noindent   \textbf{Question 4.}

\begin{parts}
  \item Explain the method of maximum likelihood estimation. Obtain the maximum likelihood estimator (MLE) of $ \theta$ based on a random sample of size $n$ from the population with probability density function:
  \[
    f(x,  \theta) = \frac{1}{ \theta} x^{(1/ \theta - 1)}, \quad 0 < x < 1, \quad  \theta > 0.
  \]\pts{7}
  \item Let $X_1, X_2, \ldots, X_n$ be a random sample from the probability density function:
  \[
    f(x,  \theta) = \frac{ \theta^x e^{- \theta}}{x!(1 - e^{- \theta})}, \quad x = 1, 2, 3, \ldots
  \]
  Find the method of moments estimator of $ \theta$. \pts{7}
\end{parts}

\medskip\rule{\linewidth}{0.2pt}

\medskip


\noindent   \textbf{Question 5.}

\begin{parts}
  \item Explain the method of construction of confidence intervals. \pts{4}
  \item Obtain the $100(1-\alpha)\%$ confidence interval for $\sigma^2$ of the normal distribution $N(\mu, \sigma^2)$ when:
  
\begin{parts}
    \item $\mu$ is known,
    \item $\mu$ is unknown.
  \end{parts}\pts{10}
\end{parts}

\medskip\rule{\linewidth}{0.2pt}

\medskip


\noindent   \textbf{Question 6.}

\begin{parts}
  \item Explain the following terms:
  
\begin{parts}
    \item Simple and composite hypothesis
    \item Type I and Type II errors
  \end{parts}\pts{7}
  \item State and prove the Neyman--Pearson fundamental lemma. \pts{7}
\end{parts}

\medskip\rule{\linewidth}{0.2pt}

\medskip


\noindent   \textbf{Question 7.}

\begin{parts}
  \item Explain the likelihood ratio test. Obtain the likelihood ratio test for testing $H_0\colon \mu = \mu_0$ against $H_1\colon \mu 
eq \mu_0$ on the basis of a random sample of size $n$ from $N(\mu, \sigma^2)$ with unknown $\sigma^2$. \pts{7}
  \item Define the loss function and risk function. On the basis of a random sample of size $n$ from a Poisson distribution with parameter $ \theta$, obtain the Bayes estimator of $ \theta$ under the squared error loss function, when the prior distribution for $ \theta$ is a Gamma distribution with parameters $a$ and $b$. \pts{7}
\end{parts}

\vfill
\rule{\linewidth}{0.4pt}

\begin{center}\small
  \href{https://bhu-exam.vercel.app/}{Click here to visit our website}\\[4pt]
  \href{https://me-aryan.vercel.app/}{Click here to visit our contributor}\\[4pt]
  \href{https://quashan-me.vercel.app/}{Click here to visit our contributor}
\end{center}

\end{document}